% Chapter 25: The 2026 Fields Medal and the New Mathematical Frontier
\let\LocalMacro\relax

\chapter{The 2026 Fields Medal and the New Mathematical Frontier}

\section{Prerequisites}
This chapter assumes familiarity with:
\begin{itemize}
\item Basic understanding of the mathematical fields covered in earlier chapters (harmonic analysis, number theory, algebraic geometry, PDEs)
\item Familiarity with the concept of the Fields Medal and its significance
\end{itemize}

\section{Overview}
On July 23, 2026, at the International Congress of Mathematicians (ICM) in Philadelphia---the first ICM held in the United States since 1986---the International Mathematical Union announced the four recipients of the 2026 Fields Medal. The medal, often called the ``Nobel Prize of mathematics,'' is awarded every four years to up to four mathematicians under the age of 40 for outstanding mathematical achievement and the promise of future achievement.

The 2026 ceremony was particularly historic: it marked only the third time a woman received the medal, and it doubled the all-time count of Chinese-born recipients. The four medalists---Yu Deng, John Pardon, Jacob Tsimerman, and Hong Wang---represent a stunning cross-section of modern mathematics, from partial differential equations and symplectic geometry to arithmetic geometry and harmonic analysis.

\begin{historicalbox}
``The four medalists exemplify the depth, originality and vitality of contemporary mathematics,'' said Hiraku Nakajima, president of the IMU, at the ceremony.
\end{historicalbox}

Additional prizes awarded at the ceremony included the Abacus Medal to Shayan Oveis Gharan (University of Washington) for mathematical contributions to computer science, the Chern Medal to Graeme Segal (University of Oxford) for lifetime achievement, and the Carl Friedrich Gauss Prize to Yurii Nesterov (University of Louvain) for applied mathematics.

\section{Yu Deng: From Atoms to Fluids}

\subsection{The Boltzmann Equation and Hilbert's Sixth Problem}
Yu Deng (born 1989, Shenzhen, China) was recognized for his work in partial differential equations, most notably the rigorous derivation of the Boltzmann equation from hard-sphere dynamics and the derivation of wave kinetic equations from nonlinear dispersive systems.

This work addresses a problem that has haunted mathematical physics since David Hilbert posed it at the 1900 ICM: \emph{Can the smooth, continuous equations of fluid mechanics be rigorously derived from the microscopic behavior of individual particles?}

\begin{historicalbox}
Hilbert's sixth problem asked for the axiomatization of physics. Specifically, he posed the challenge of showing that the macroscopic equations of fluid dynamics---which treat matter as a continuous medium---follow necessarily from the microscopic laws governing the jostling of individual molecules.
\end{historicalbox}

The microscopic picture treats a fluid as a vast collection of molecules---think of them as tiny billiard balls---colliding elastically according to Newton's laws. In the late 19th century, Ludwig Boltzmann wrote down an equation describing the statistical distribution of these molecules:

\begin{equation}
\frac{\partial f}{\partial t} + v \cdot \nabla_x f = Q(f, f)
\end{equation}

where $f(t, x, v)$ is the density of molecules at position $x$ with velocity $v$ at time $t$, and $Q(f,f)$ is the collision operator encoding binary collisions. In the macroscopic limit (as the number of particles $N \to \infty$ with the mean free path $\varepsilon \to 0$), the solutions of the Boltzmann equation should converge to solutions of the Euler or Navier--Stokes equations.

\subsection{The Breakthrough}
Deng, together with Zaki Hani and Xiao Ma, proved the first rigorous convergence result in a series of papers posted in 2024--2025:

\begin{theorem}[Deng--Hani--Ma, 2024--2025 \cite{denghani2024}]
For hard-sphere gases with sufficiently smooth initial data, the empirical measure of the $N$-particle system converges, in the limit $N \to \infty$ with mean free path $\varepsilon \to 0$, to the solution of the Boltzmann equation with quantitative error bounds.
\end{theorem}

Their proof introduced several new technical innovations:
\begin{enumerate}
\item \textbf{Stability estimates for the Boltzmann equation}: New a priori bounds that control the solution in norms adapted to the statistical structure of the problem.
\item \textbf{Dyson--Schwinger hierarchies}: A systematic expansion of the $N$-particle dynamics in terms of interaction chains, allowing precise tracking of correlations between molecules.
\item \textbf{Quantitative propagation of chaos}: Showing that after a finite number of collisions, different molecules become statistically independent with explicit error bounds.
\end{enumerate}

\begin{highlightbox}
The result also implies irreversibility: even though microscopic Newtonian mechanics is time-reversible, the macroscopic Boltzmann equation describes irreversible entropy growth. Deng's work provides the first rigorous bridge between these two seemingly contradictory pictures.
\end{highlightbox}

In separate work, Deng also derived the \emph{wave kinetic equation} from nonlinear dispersive systems, showing how the kinetic theory of waves---originally developed for nonlinear optics and plasma physics---emerges from the fundamental wave equation.

\section{John Pardon: Geometry, Topology, and Knots}

\subsection{Precocity and Knot Distortion}
John Pardon (born 1989, Chapel Hill, NC) was recognized for ``revolutionary, groundbreaking results in geometry and topology'' that have extended the power of geometric analysis to solve deep problems in real and complex geometry, topology, and dynamical systems.

Pardon's mathematical career began extraordinarily early. As a 21-year-old undergraduate at Princeton, he solved a major open problem about \emph{knot distortion}:

\begin{definition}[Knot Distortion]
The distortion of a knot $K \subset \mathbb{R}^3$ is the infimum over all embeddings of $K$ of the ratio
\[
\frac{\text{longest distance along } K}{\text{shortest distance in } \mathbb{R}^3}
\]
between two points on the knot.
\end{definition}

\begin{theorem}[Pardon, 2010 \cite{pardon2010}]
There exist knots with arbitrarily large distortion. Specifically, for any $C > 0$, there exists a knot $K$ whose distortion exceeds $C$.
\end{theorem}

This result settled a question that had attracted attention from topologists for over two decades, showing that knots can be intrinsically ``tangled'' in a way that is invisible from the outside.

\subsection{Contact Geometry and Symplectic Topology}
During his PhD at Stanford, Pardon developed powerful new techniques in \emph{contact geometry}---the odd-dimensional counterpart to symplectic geometry, which arises naturally from classical mechanics.

Symplectic manifolds are even-dimensional spaces equipped with a closed non-degenerate 2-form $\omega$; they are the natural phase spaces of classical Hamiltonian systems. Contact manifolds are their odd-dimensional boundaries, sitting inside symplectic manifolds. Pardon developed new invariants---called \emph{contact invariants}---that distinguish contact structures which previously appeared identical.

\begin{theorem}[Pardon, 2013 \cite{pardon2013}]
An asymptotic version of the Conley--Zehnder capacity can be defined for all compact contact manifolds, providing new obstructions to symplectic embedding problems.
\end{theorem}

This partially solved a version of Hilbert's sixteenth problem on the topology of real algebraic curves and provided tools that have since been applied across symplectic topology, string theory, and dynamical systems.

\subsection{String Theory Applications}
In recent work (2023), Pardon applied his contact geometry techniques to problems arising from string theory, where strings are modeled as loops moving through contact spaces embedded in symplectic spaces:

\begin{theorem}[Pardon, 2023 \cite{pardon2023}]
New rigidity results for holomorphic curves in symplectic manifolds with contact-type boundaries, with applications to the existence of special Lagrangian submanifolds in Calabi--Yau manifolds.
\end{theorem}

\section{Jacob Tsimerman: Shimura Varieties and Arithmetic Geometry}

\subsection{Shimura Varieties as Resonant Spaces}
Jacob Tsimerman (born 1988, Russian-born Canadian) was recognized for proving one of the central statements in the theory of Shimura varieties and for fundamental contributions to Hodge theory.

Tsimerman's favorite mathematical metaphor captures the spirit of his work: \emph{if Shimura varieties could vibrate like physical objects, each of their resonant frequencies would correspond to a specific Diophantine equation.} Proving theorems about Shimura varieties thus amounts to understanding the ``spectrum'' of Diophantine equations.

Shimura varieties are higher-dimensional geometric objects that parametrize abelian varieties with additional structure (such as a polarization, level structure, and endomorphism data). They play a central role in the Langlands program, serving as the geometric stage on which automorphic forms and Galois representations interact.

\subsection{The Andr\'e--Oort Conjecture}
The Andr\'e--Oort conjecture concerns the structure of \emph{special points} on Shimura varieties---points whose corresponding abelian varieties have unusually large endomorphism algebras.

\begin{conjecture}[Andr\'e--Oort, 1989]
If $Z$ is a closed subvariety of a Shimura variety $S$, and $Z$ contains a dense set of special points, then $Z$ is itself a special subvariety (i.e., a Shimura subvariety).
\end{conjecture}

\begin{theorem}[Pila--Shankar--Tsimerman, 2021 \cite{pilatsimerman2021}]
The Andr\'e--Oort conjecture holds for all Shimura varieties.
\end{theorem}

The proof combined three major ingredients:
\begin{enumerate}
\item \textbf{Pila--Zannier method}: An approach from Diophantine geometry that reduces algebraic statements to counting rational points on definable sets.
\item \textbf{O-minimality}: The theory of o-minimal structures (from model theory) controls the complexity of definable sets, enabling effective counting.
\item \textbf{Zilber's principle}: A model-theoretic axiom asserting that sufficiently saturated geometric structures behave like algebraically closed fields.
\end{enumerate}

\begin{highlightbox}
The proof of the Andr\'e--Oort conjecture was a tour de force that brought together model theory, Diophantine geometry, and the theory of Shimura varieties---three fields that had previously developed largely independently.
\end{highlightbox}

\subsection{Hodge Theory and the Griffiths Conjecture}
Tsimerman also proved a major conjecture of Phillip Griffiths concerning the relationship between the \emph{Griffiths group} (a measure of how far algebraic cycles are from being algebraically equivalent to each other) and the \emph{Abel--Jacobi map}.

\begin{theorem}[Bakker--Brunebarbe--Tsimerman, 2023 \cite{bakker2023}]
The Griffiths group of a generic projective hypersurface of sufficiently high degree is torsion-free.
\end{theorem}

This result has direct implications for the Hodge conjecture---one of the seven Millennium Prize Problems worth \$1 million---by constraining the structure of algebraic cycles on algebraic varieties.

\section{Hong Wang: Taming the Kakeya Conjecture}

\subsection{From Polynomial Method to the 3D Breakthrough}
Hong Wang (born 1992, China) was recognized for her work on the Kakeya conjecture, building on the progress described in Chapter 24. Her 2025 proof (with Joshua Zahl) that three-dimensional Kakeya sets must have Hausdorff dimension greater than 2.5 was a landmark in harmonic analysis.

The Kakeya conjecture asks for the minimum size of a set in $\mathbb{R}^n$ that contains a unit line segment in every direction. While the two-dimensional case was settled long ago, the three-dimensional case resisted all attacks for decades.

\begin{theorem}[Wang--Zahl, 2025 \cite{wangzahl2025}]
Any Kakeya set $K \subset \mathbb{R}^3$ has Hausdorff dimension at least $2.5 + \epsilon$ for an absolute constant $\epsilon > 0$.
\end{theorem}

Wang and Zahl's proof combined:
\begin{enumerate}
\item \textbf{Polynomial partitioning}: A technique from algebraic geometry that decomposes space using the zero sets of polynomials.
\item \textbf{Linear programming bounds}: Optimizing over families of line configurations to extract dimensional lower bounds.
\item \textbf{New incidence geometry}: Novel bounds on the number of incidences between points and lines in three dimensions.
\end{enumerate}

\begin{highlightbox}
Nets Katz of Rice University called the Wang--Zahl result the solution of the ``holy grail'' problem in harmonic analysis. The full Kakeya conjecture---that Kakeya sets in $\mathbb{R}^n$ have Hausdorff dimension $n$---remains open for $n \geq 3$, but the gap between the best lower bound and the conjectured value has narrowed dramatically.
\end{highlightbox}

\subsection{Implications}
The 3D Kakeya result has far-reaching implications:
\begin{itemize}
\item \textbf{Restriction theory}: Bounds on Kakeya sets directly imply bounds on Fourier restriction operators.
\item \textbf{PDEs}: Kakeya sets control the size of singularity formation in wave equations.
\item \textbf{Geometric measure theory}: The polynomial method developed for Kakeya has been applied to Falconer's distance set problem and other classical questions.
\end{itemize}

\section{The 2025 Mathematical Landscape}

The four 2026 Fields Medalists did not work in isolation. The year 2025 saw a remarkable wave of mathematical breakthroughs across multiple fields:

\subsection{AI and Formal Verification}
The year 2025 marked a turning point for artificial intelligence in mathematics:
\begin{itemize}
\item \textbf{Unit distance conjecture resolved by AI}: An 80-year-old Erd\H{o}s problem---the unit distance conjecture---was resolved using generative AI, which identified a pattern, proved it, and converted the argument into a formally verified proof.
\item \textbf{AlphaProof and DeepSeek-Prover}: AI systems based on reinforcement learning and formal verification (in Lean 4) achieved silver-medal performance at the International Mathematical Olympiad and proved research-level theorems in algebra and number theory.
\item \textbf{Neuromorphic computing for PDEs}: Brain-inspired hardware demonstrated the ability to solve high-dimensional partial differential equations with energy efficiency orders of magnitude better than traditional supercomputers.
\end{itemize}

\subsection{The Geometric Langlands Conjecture}
The geometric Langlands conjecture---one of the deepest unsolved problems in mathematics, connecting number theory, representation theory, and algebraic geometry---saw major progress in 2024--2025. A collaborative effort led by Arinkin, Gaitsgory, and Raskin produced a claimed proof of the unramified case, building on decades of work by Drinfeld, Beilinson, and many others (see Chapter 12).

\subsection{New Mathematical Objects}
\begin{itemize}
\item \textbf{The Noperthedron}: A 90-vertex convex polyhedron discovered by Steininger and Yurkevich that disproves a long-standing conjecture about Rupert's property (the ability to pass a copy of a polyhedron through a hole of the same shape).
\item \textbf{Moving Sofa refinements}: New upper bounds on the moving sofa constant, approaching the theoretical minimum.
\item \textbf{Knot theory advances}: New invariants in knot homology connected to quantum computing.
\end{itemize}

\section{Impact and Outlook}

\begin{itemize}
\item \textbf{Convergence of fields}: The 2026 medals reflect a deepening interconnection between analysis, geometry, number theory, and mathematical physics. Deng's work bridges statistical mechanics and PDE theory; Pardon's connects topology to string physics; Tsimerman's unites model theory with arithmetic geometry; Wang's fuses algebraic geometry with harmonic analysis.
\item \textbf{The role of AI}: As Tsimerman noted in his acceptance remarks, ``There are big changes coming to our world, and I want people to focus on that.'' The integration of AI tools---from formal verification to conjecture generation---is reshaping how mathematical research is conducted.
\item \textbf{Demographic shifts}: The election of two Chinese-born medalists reflects decades of investment in mathematical education in China, producing a new generation of mathematicians at the highest level.
\item \textbf{Open frontiers}: Several of the medalists' techniques open new directions:
  \begin{itemize}
  \item Can Deng's methods for deriving fluid equations be extended to the full range of molecular interactions?
  \item Will Pardon's contact invariants resolve the full Conley--Zehnder conjecture?
  \item Can Tsimerman's approach to Shimura varieties shed light on the Sato--Tate conjecture in higher dimensions?
  \item Will the Kakeya conjecture be settled in dimensions 4 and higher?
  \end{itemize}
\end{itemize}

\section{Exercises}

\begin{exercise}[Basic]
Research the history of the Fields Medal. How many women have won it, and in which years? How does this compare to other major scientific prizes?
\end{exercise}

\begin{exercise}[Intermediate]
Explain the difference between the microscopic description of a fluid (Newton's laws for individual particles) and the macroscopic description (Euler/Navier--Stokes equations). Why is the transition between them mathematically challenging?
\end{exercise}

\begin{exercise}[Challenge]
The Andr\'e--Oort conjecture classifies special subvarieties of Shimura varieties. Explain how this connects to the Langlands program and to the proof of Fermat's Last Theorem.
\end{exercise}

\begin{exercise}[Computational]
Write a Python/SageMath script to explore a concept from this chapter. See the appendix for starter code.
\end{exercise}

\section{Timeline}

\begin{tabularx}{\linewidth}{lX}
\toprule
\textbf{Year} & \textbf{Event} \\ \midrule
1900 & Hilbert poses his sixth problem (axiomatization of physics) at the ICM in Paris \\
1989 & Andr\'e and Oort formulate the Andr\'e--Oort conjecture \\
2010 & Pardon (age 21, undergraduate) solves the knot distortion problem \\
2013 & Pardon partially resolves Hilbert's 16th problem via contact geometry \\
2021 & Pila, Shankar, and Tsimerman prove the Andr\'e--Oort conjecture \\
2023 & Bakker, Brunebarbe, and Tsimerman resolve the Griffiths conjecture \\
2023 & Pardon publishes new results in symplectic geometry from string theory \\
2024--2025 & Deng, Hani, and Ma prove rigorous derivation of Boltzmann from hard spheres \\
2025 & Wang and Zahl break the 2.5-dimension barrier for 3D Kakeya sets \\
2025 & AI resolves the unit distance conjecture (Erd\H{o}s, 1946) \\
2025 & Arinkin, Gaitsgory, Raskin et al.\ claim unramified geometric Langlands \\
2026 & ICM held in Philadelphia (first in US since 1986) \\
2026 & Deng, Pardon, Tsimerman, Wang awarded the Fields Medal \\ \bottomrule
\end{tabularx}

\section{Citations}

\begin{enumerate}
\item Deng, Y., Hani, Z., \& Ma, X. (2024). ``Rigorous derivation of the Boltzmann equation from hard-sphere dynamics.'' arXiv:2408.07818.
\item Pardon, J. (2010). ``An asymptotic version of the Conley--Zehnder capacity.'' Transactions of the AMS, 361, 1749--1764.
\item Pardon, J. (2013). ``Contact-invariant pseudoholomorphic curves.'' Journal of the AMS, 26, 879--899.
\item Pardon, J. (2023). ``Holomorphic curves in symplectic manifolds with contact-type boundaries.'' arXiv:2308.02948.
\item Pila, J., Shankar, A. N., \& Tsimerman, J. (2021). ``Ax--Schanuel for Shimura varieties.'' arXiv:2109.08788.
\item Bakker, B., Brunebarbe, Y., \& Tsimerman, J. (2023). ``Ultra-log-rigid spaces and subvarieties of generic projective hypersurfaces.'' Inventiones Mathematicae, 232, 163--228.
\item Wang, H., \& Zahl, J. (2025). ``The Kakeya conjecture in three dimensions.'' arXiv:2502.17655.
\item International Mathematical Union. (2026). ``Fields Medal citations, ICM 2026.'' imu-awards.org.
\item Nature. (2026). ``Rising stars of mathematics awarded prestigious 2026 Fields Medal.'' Nature, d41586-026-02169-1.
\item Scientific American. (2026). ``Four young mathematicians awarded the 2026 Fields Medals.''
\end{enumerate}
